\section{Heat-pump performance calculation}
\label{app:cop}

The heat-pump COPs are calculated exogenously using an adaptation of the
central-heating COP approximation implemented in PyPSA-Eur~\cite{pypsaeur2025}.
The approximation is based on Jensen et al.~\cite{jensen2018cop}, using the
parameterisation evaluated by Pieper et al.~\cite{pieper2020cop}.
For each $i\in\hpset$, the logarithmic mean absolute temperatures (in kelvin) are
\begin{equation}
\widetilde T_{j,t}^i
=\frac{T_{j,t}^{i,\mathrm{hot}}-T_{j,t}^{i,\mathrm{cold}}}
  {\ln\!\big(T_{j,t}^{i,\mathrm{hot}}/T_{j,t}^{i,\mathrm{cold}}\big)},
\label{eq:log_mean_temperature}
\end{equation}
where $j\in\{\mathrm{sink},\mathrm{src}\}$ indexes the heat sink and source.
The corresponding ideal Lorenz COP is
\begin{equation}
\mathrm{COP}_t^{\mathrm{L},i}
=\frac{\widetilde T_{\mathrm{sink},t}^i}
  {\widetilde T_{\mathrm{sink},t}^i-\widetilde T_{\mathrm{src},t}^i}.
\label{eq:lorenz_cop}
\end{equation}
It is adjusted for refrigerant properties, heat-exchanger pinch temperature,
and compressor efficiency:
\begin{equation}
\begin{aligned}
\mathrm{COP}_t^i={}&\mathrm{COP}_t^{\mathrm{L},i}F_t^i
  \eta^{\mathrm{is}}(1-r_t^i)\\
 &+(1-\eta^{\mathrm{is}})-l^{\mathrm{HP}}.
\end{aligned}
\label{eq:corrected_cop}
\end{equation}
Here, $F_t^i$ and $r_t^i$ are the correction factors of the underlying
approximation. Ammonia is assumed as refrigerant, with an isentropic compressor
efficiency of 0.75, a pinch-point temperature difference of 2~K, and no
additional heat-loss term. COP values below one or temperature lifts below
5~K are excluded~\cite{zajacs2020lowtemperature}.

District-heating return and supply temperatures are fixed at
$40^\circ\mathrm{C}$ and $75^\circ\mathrm{C}$, respectively.
The assumed source inlet/outlet temperatures are $12/8^\circ\mathrm{C}$ for
City2 and $10.5/4^\circ\mathrm{C}$ for the groundwater heat pump.
These assumptions give constant COPs of 3.51 and 3.30, respectively.
